\section{Equivariant quotients}
\label{sec:equivariant-quotients}

An equivalence relation compatible with an action allows renaming to pass
to equivalence classes. The construction concerns the selected action and
the relation; it requires no support hypothesis. We first state it for
monoids, then specialize to finite permutations.

\subsection{Invariant relations and the induced action}

Let a monoid $M$ act on $X$, and let $\sim$ be an equivalence relation on
$X$ satisfying
\[
  x\sim y\quad\Longrightarrow\quad m\act x\sim m\act y
  \qquad(m\in M).
\]
Only preservation is needed. For a group action, applying the same
condition to $m^{-1}$ also gives reflection. Write $q(x)=[x]$ for the
ordinary quotient projection.

\begin{proposition}[Canonical quotient action]
\label{prop:quotient-action}
The equation
\[
  m\act q(x)=q(m\act x)
\]
defines an action of $M$ on $X/{\sim}$. The projection is surjective and
commutes with the actions.
\end{proposition}

\begin{proof}
If $x\sim y$, compatibility gives $m\act x\sim m\act y$, so the
formula is independent of the representative. On a class $q(x)$ the
action laws reduce to those on $X$:
\[
  1\act q(x)=q(x),\qquad
  (mn)\act q(x)=q(m\act(n\act x))=m\act(n\act q(x)).
\]
Every class has a representative, which proves both that these laws hold
on the whole quotient and that $q$ is surjective. The defining equation,
read in the opposite direction, is the action-preservation law for $q$.
\end{proof}

The group case is the quotient $G$-set construction of
Pitts~\cite[Section~1.7, equation~(1.43)]{pitts2013}. No nonempty-carrier
assumption is used: a quotient of an empty carrier is empty, and its
action laws are still meaningful. Neither decidability of the relation
nor any assumptions about atoms enter the construction.

Lean uses \texttt{Setoid X} for the equivalence relation and its laws.
The carrier is \texttt{Quotient s}.
Compatibility with the selected scalar operation is the ordinary
proposition \texttt{SMulInvariant M s}.
Quotient mapping in Mathlib descends each operation
$x\mapsto m\act x$, with representative computation given by
\path{Quotient.map_mk}. This part needs only \texttt{SMul M X},
without monoid laws. The resulting projection can be packaged using
Mathlib's existing \texttt{MulActionHom}; it remains an ordinary function
under coercion.

The monoid laws transfer along the projection using Mathlib's
\path{Function.Surjective.mulAction}.
That result presupposes a scalar operation on the
target; quotient mapping supplies it. The inherited operation is
definitionally the same one, so the public representative equation
\path{QuotientAction.smul_mk} applies to both interfaces. The construction
uses quotient elimination, without selecting representatives as data, and
keeps the scalar and carrier universes independent.

The action is installed explicitly by a client. Compatibility is an ordinary
proof argument, not an instruction for typeclass search to select a new
action on every quotient. Consequently the projection laws refer to this
constructed action even if another action is available on the same
underlying quotient type. A different action must establish its own
projection law. This preserves the selected-action distinction of
Section~\ref{sec:actions} while making canonical quotient computation
available through a small public interface.

\subsection{Support and nominality transfer}

Now specialize to selected actions of $\Perm(A)$, with $A$ arbitrary.
Equivariant maps preserve sufficient support bounds by
Proposition~\ref{prop:support-basic}. Surjectivity gives a further consequence
for whole carriers.

\begin{proposition}[Nominality under equivariant surjections]
\label{prop:surjective-nominality}
Let $f:X\to Y$ be an equivariant surjection between selected permutation
actions. If the action on $X$ is nominal, so is the action on $Y$.
Neither infinitude of $A$ nor nonemptiness of $X$ is required.
\end{proposition}

\begin{proof}
For $y\in Y$, surjectivity supplies $x\in X$ with $f(x)=y$.
Nominality of $X$ supplies a finite supporting bound for $x$; equivariance
makes the same bound support $y$. The preimage is used only to prove
existence of a finite bound, not to define an operation on $Y$.
\end{proof}

This is \path{Equivariant.nominal_of_surjective}. Its target action is
already selected; the theorem supplies only the proof-only nominality
certificate. It is an explicit theorem, not an instance-search rule that
could choose an action through an arbitrary surjection.

\begin{proposition}[Canonical quotient support]
\label{prop:quotient-support}
Suppose $\sim$ is invariant under the selected action on $X$, and give
$X/{\sim}$ its canonical action. Every finite set supporting $x$ also
supports $q(x)$. Thus finite supportedness passes to individual classes,
and a nominal source has a nominal quotient. If $A$ is infinite and $x$
is finitely supported, then
\[
  \operatorname{supp}_A(q(x))\subseteq\operatorname{supp}_A(x).
\]
The individual conclusions require no nominality of the whole source.
\end{proposition}

\begin{proof}
The projection is equivariant by Proposition~\ref{prop:quotient-action}.
Apply the existing image support theorem to any supplied finite bound.
Apply Proposition~\ref{prop:surjective-nominality} for the carrier statement.
Under infinite atoms, apply image minimality to the source's least supporting
bound to obtain the inclusion.
\end{proof}

The Lean interface keeps the same separation as the preceding support
calculus. A source certificate \texttt{hx} maps to a certificate for its
class, and \path{FinitelySupported.support_map_subset} gives the elementwise
inclusion. On a nominal source, the canonical quotient certificate permits
the convenient carrier statement \path{QuotientAction.support_mk_subset}.
These are applications of the established map laws; no new least-support
choice or support-intersection argument is needed. In particular, a quotient
of a nominal action over finite atoms is still nominal, even though the
general least-support interface requires infinite atoms.

The canonical-action qualification is essential. For the equality relation
on atoms, the canonical quotient action still renames the represented atom.
One can instead give the same quotient carrier a trivial action. If $a\ne b$,
this other action fixes $q(a)$, while
$q(\swap a b\act a)=q(b)\ne q(a)$. The projection is then not equivariant.
A certificate or equation about the canonical action cannot be applied to
this other action solely because their carrier types agree.

\subsection{Strict decrease and supported fibers}

\begin{example}[Forgetting one component]
\label{ex:quotient-strict-support}
Give $A\times A$ the componentwise action and relate $(a,b)$ to $(a',b')$
exactly when $a=a'$. This relation is invariant. Over infinite atoms, its
canonical quotient satisfies
\[
  \operatorname{supp}_A(q(a,b))=\{a\},\qquad
  \operatorname{supp}_A(a,b)=\{a,b\}.
\]
For $a\ne b$, the quotient support is a strict subset of the representative's
support: it retains dependence on $a$ and loses dependence on $b$.
\end{example}

\begin{proof}
The classes $q(a,b)$ and $q(a,a)$ agree. The latter has support contained
in $\{a\}$ by the quotient bound and the exact atom/product formulas.
Conversely, first projection descends to an equivariant function
$(A\times A)/{\sim}\to A$ sending $q(a,b)$ to $a$. Image support inclusion for this
function puts $\{a\}$ inside the class's support, giving equality.
The pair formula then proves strictness when $a\ne b$.
\end{proof}

There is nevertheless an exact characterization using all supported
preimages. It applies beyond quotients, without nominality of either
whole carrier or surjectivity of the map.

\begin{proposition}[Support in a supported fiber]
\label{prop:supported-fiber}
Assume $A$ is infinite, $f:X\to Y$ is equivariant, and $x\in X$ is
finitely supported. Let
\[
  F_{f,x}=\{z\in X\mid z\text{ is finitely supported and }f(z)=f(x)\}.
\]
Then
\[
  \operatorname{supp}_A(f(x))
  =\bigcap_{z\in F_{f,x}}\operatorname{supp}_A(z).
\]
\end{proposition}

\begin{proof}
The fiber is nonempty, since it contains $x$. Image support inclusion
puts the support of $f(x)$ inside every support on the right.
For the converse, suppose $a\notin\operatorname{supp}_A(f(x))$.
Choose $b\notin\{a\}\cup\operatorname{supp}_A(x)$; image inclusion also
puts $b$ outside the image support. Hence $\tau=\swap a b$ fixes $f(x)$.
Equivariance gives $f(\tau\act x)=f(x)$, and renaming preserves finite
supportedness, so $\tau\act x\in F_{f,x}$. Support transport gives
\[
  \operatorname{supp}_A(\tau\act x)
  =\tau[\operatorname{supp}_A(x)].
\]
This set omits $a$, since its inverse image under $\tau$ is the excluded
atom $b$. Thus $a$ is also absent from the intersection.
\end{proof}

The theorem \path{FinitelySupported.mem_support_map_iff} uses
membership: an atom belongs to image support exactly
when it belongs to the support of every finitely supported preimage.
Unsupported values receive no artificial least support. The proof uses
the established image bound, outside-swap fixation and transport law;
it reconstructs none of them. Infinite atoms provide the single fresh
choice by finite avoidance; countability is unnecessary. Classical equality
for the swap and finite image remains inside the proof.

\begin{proposition}[Support of a quotient class]
\label{prop:quotient-support-intersection}
Assume $A$ is infinite and $X$ is nominal. For an invariant equivalence
relation and any class $c$ under the canonical action,
\[
  \operatorname{supp}_A(c)
  =\bigcap_{\{x\in X\mid q(x)=c\}}\operatorname{supp}_A(x).
\]
\end{proposition}

Apply Proposition~\ref{prop:supported-fiber} to a representative of $c$;
nominality supplies support for all representatives. This is the quotient
characterization of Pitts~\cite[Proposition~2.30, pp.~45--46]{pitts2013}.
Its fresh-swap proof also explains the weaker hypotheses of the general
fiber result. The Lean theorem \path{QuotientAction.support_eq_iInter}
uses quotient induction and set extensionality, indexing the intersection
by representatives together with their projection equality. The right side
is an intersection in \texttt{Set A}, whereas each individual support is
a Finset. This is an intersection of supports of different objects, not an
assumption that arbitrary intersections of supporting bounds support one
fixed object.

\begin{example}[No representative need attain class support]
\label{ex:quotient-support-nonattainment}
Over an infinite atom carrier, identify every atom with every other atom.
The resulting class is fixed
by all permutations, so its least support is empty. Every representative
still has singleton support. Thus the intersection formula does not give
a representative whose support equals the class support.
\end{example}

Nor does support of a class imply support of a representative. For example,
the identity in the pointwise function action on $\mathbb N\to\mathbb N$
is unsupported, as shown in Section~\ref{sec:least-support}. In the universal
quotient of that carrier its class is invariant and therefore empty-supported.
The fiber characterization quantifies only over supported preimages; it
does not reverse finite supportedness or carrier nominality through a
surjection.

\subsection{Supported contexts and freshness}

Assume $A$ is infinite. The established general freshness laws immediately
give
\[
  c\mathrel{\#}x\ \Longrightarrow\ c\mathrel{\#}q(x),\qquad
  x\mathrel{\#}c\ \Longrightarrow\ q(x)\mathrel{\#}c.
\]
Indeed, quotienting can only shrink the relevant support, preserving its
disjointness from the context's support. The context may have a different
carrier and universe. Individual support certificates for $c$ and $x$
suffice; their entire carriers need not be nominal. Nested products and
finite atom sets use the decomposition laws of Section~\ref{sec:freshness}.
For example, freshness of $(c,(d,e))$ for $x$ gives freshness of each
component for $q(x)$.
The implication need not reflect freshness: in
Example~\ref{ex:quotient-strict-support}, $b\mathrel{\#}q(a,b)$ for
$a\ne b$, while $b$ is not fresh for the original pair.

These consequences use the existing freshness-map theorems with the
projection. They concern supports of objects. The support-reflection statement
for predicate pullback in Proposition~\ref{prop:predicate-pullback} instead
compares dependence of $P$ with that of $P\circ q$; it is a distinct theorem
about predicates, not an equality between the supports of an object and its
representatives.
